% ---------------------------------------------------------------------------
\section{Feature Engineering}
\label{sec:features}
% ---------------------------------------------------------------------------

\subsection{Derived features and the log target}

Eighteen derived variables are computed from the cleaned table. Every one is
row-wise, so the identical transformation applies to a single listing at inference
time and no train--serve skew is possible. Table~\ref{tab:engineered} groups them by
business purpose. The composites exist because several raw attributes measure the
same underlying quantity: usable area is the sum of basement, first-floor and
second-floor area, and bathroom provision is better summarised by a weighted count
that treats a half bathroom as half a full one. Age is expressed as age at sale and
years since remodelling rather than as construction years, because an agent reasons
about how old a house is, not about the calendar year in which it was built.

\begin{table}[htbp]
\centering
\caption{Engineered features by business purpose. The first sixteen are model inputs;
the last two are analysis variables and are excluded from the model matrix.}
\label{tab:engineered}
\small
\begin{tabularx}{\linewidth}{@{}L{0.18\linewidth}L{0.36\linewidth}X@{}}
\toprule
\textbf{Family} & \textbf{Features} & \textbf{Rationale} \\
\midrule
Floor area & \f{total\_sf}, \f{total\_porch\_sf} & Usable covered area; aggregate outdoor amenity space \\
Bathrooms & \f{total\_bath} & Full-bathroom equivalents in one count \\
Age and renovation & \f{age\_at\_sale}, \f{remod\_age}, \f{is\_renovated}, \f{renovation\_scope\_code} & Depreciation and vintage; depth of the last permitted renovation \\
Presence flags & \f{has\_pool}, \f{has\_garage}, \f{has\_fireplace}, \f{has\_second\_floor} & Separates absence of an amenity from a small amount of it \\
Ratios & \f{liv\_area\_ratio}, \f{garage\_area\_ratio}, \f{rooms\_per\_1000sf}, \f{avg\_room\_size} & Development intensity; layout spaciousness \\
Quality & \f{quality\_condition\_index} & Joint effect of quality and maintenance condition \\
Analysis only & \f{price\_per\_sqm}, \f{log\_price} & Reporting metric and target; both use the price and are never inputs \\
\bottomrule
\end{tabularx}
\end{table}

The model is trained on $\log(1 + \text{SalePrice})$. The transformation reduces
skewness from 1.881 to 0.122, it matches the multiplicative structure of housing
prices, and it stabilises the variance of expensive observations so that a 700,000
\USD{} house does not dominate the fit. It also has a consequence that is often
overlooked: because the exponential of a mean log price is not the mean price,
back-transforming with $\exp(\hat{\mu}) - 1$ estimates a conditional median and is
biased low by roughly half the residual variance unless a smearing correction is
applied \citep{hastie2009esl}. No smearing correction is applied here. The
consequence is visible in the diagnostics: the mean signed error is $-853$ \USD, so
the model predicts slightly below the realised price on average. Given that
over-valuation is weighted more heavily in Equation~\eqref{eq:loss}, the direction is
commercially convenient, but the report presents it as a bias rather than as
accuracy.

\subsection{Encoding, selection and collinearity}

Numerical inputs are median-imputed and standardised; categorical inputs are imputed
with the most frequent level and one-hot encoded with a minimum frequency of ten
observations, so rare levels are grouped instead of producing indicators estimated on
a handful of rows. Unseen levels at inference time are ignored rather than treated as
errors. Tree learners use either the same one-hot matrix without scaling or
LightGBM's native categorical handling, so preprocessing is part of the fitted
pipeline and is identical offline and online. After encoding the matrix contains 102
inputs: 40 categorical, 29 structural, 12 structural composites, 10 neighbourhood, 5
macro-financial, 2 locational and 4 age or renovation variables.

Selection proceeds in three steps. Four uninformative columns are removed because
they are constant or almost constant in the sample: street type, utilities, pool
quality and the raw renovation flag, which duplicates the engineered
\f{is\_renovated} flag. A second step applies the leakage policy described below.
The third measures multicollinearity on the 55-feature identifiable numeric core, and
this step repays attention. The deliberate algebraic composites are excluded from the
computation, because including a composite together with its components would make
the design rank-deficient and produce infinite factors that say nothing about the
data. The remaining core is pruned iteratively: the feature with the largest factor
above ten is removed, factors are recomputed, and the process repeats until none
exceeds the threshold. Eight features are removed and 47 are retained, with a maximum
retained factor of 8.95 for garage area as a share of lot area. The eight removed features begin with the amenity score
(351.1), which is almost a linear combination of the distances, ratings and crime
variables from which it was composed, and continue with \f{has\_pool} (62.7),
\f{mortgage\_rate\_pct} (57.3), \f{second\_flr\_sf} (42.6), \f{tot\_rms\_abv\_grd}
(41.5), \f{local\_unemployment\_pct} (40.8), \f{quality\_condition\_index} (39.8)
and \f{renovation\_scope\_code} (13.4).



The pruning is a diagnostic of redundancy, not a filter applied to the model matrix.
Features such as the number of rooms above grade and the quality--condition index
remain inputs, and the regularised linear model handles their collinearity by
shrinking and distributing coefficient mass across correlated variables. The
practical implication is interpretive: individual coefficients on collinear features
must not be read as independent effects, which is why the report relies on grouped
SHAP shares rather than on raw coefficient magnitudes.

The leakage exclusion list is enforced in configuration and asserted by the test
suite. It removes direct functions of the target (\f{SalePrice}, \f{log\_price},
\f{price\_per\_sqm}); variables recorded only after the transaction closes
(\f{sale\_condition}, \f{sale\_type}); marketing outcomes (\f{days\_on\_market},
\f{listing\_date}); sale timing replaced by the macro context (\f{sale\_year},
\f{sale\_month}, \f{yr\_sold}, \f{mo\_sold}); and identifiers. The most consequential
exclusion is days on market. It is plausibly predictive, with a Spearman correlation
of $-0.435$ with price, but it is observed only after a listing has sold and is
therefore unavailable when the estimate is produced.
